\section{Experimentelles Setup und Methodik}

\subsection{Simulationsumgebung}

Für die praktische Analyse wurde eine lokale Loopback-Umgebung (127.0.0.1) auf einem
Windows-System eingerichtet. OpenSSL 3.6.1 diente auf Client und Server als kryptographische Grundlage~\cite{openssl}. Der Datenverkehr wurde parallel mit Wireshark aufgezeichnet, um den Handshake-Verlauf und Fehlermeldungen auf Paketebene zu analysieren.

\subsection{TLS-Fehlersimulation und Log-Analyse}

Es werden ein Downgrade- und ein Handshake-Fehlerfall simuliert, um zu untersuchen, wo Schutzmechanismen greifen oder versagen.

% ── Szenario 1 ────────────────────────────────────────────────────────────────
\subsection{Szenario~1: Nicht-erfolgreicher Downgrade (Handshake-Fehlerfall)}

Im ersten Experiment wurde untersucht, wie ein moderner Client reagiert, wenn ein Server
ein veraltetes, unsicheres Protokoll erzwingt. Dies simuliert das Verhalten bei einer
serverseitigen Fehlkonfiguration oder einem Manipulationsversuch im
Netzpfad~\cite{sentinelone}.

\subsection*{Vorbereitung der Zertifikate}

Zunächst wurde durch die Nutzung von OpenSSL ein selbstsigniertes X.509-Zertifikat mit
einem 2048-Bit RSA-Schlüsselpaar generiert. Der Parameter \texttt{-nodes} (No-DES) sorgt
dafür, dass der private Schlüssel unverschlüsselt bleibt, um einen reibungslosen
Serverstart ohne Passwortabfrage zu ermöglichen~\cite{openssl}:

\begin{lstlisting}[style=consolestyle]
openssl req -x509 -newkey rsa:2048 -keyout key.pem -out cert.pem \
-days 365 -nodes
\end{lstlisting}

\subsection*{Verbindungsaufbau}

Der simulierte Server wurde so gestartet, dass er ausschließlich das veraltete Protokoll
TLS~1.0 anbietet. Anschließend versuchte der reguläre Client, eine sichere Verbindung
zu diesem Port aufzubauen. Siehe Anhang~\ref{app:sz1-befehle}.

\subsection*{Analyse der Handshake-Schritte}

Der Verbindungsaufbau wird durch den Client gestartet, worauf der Server antwortet und
versucht, TLS~1.0 aufzuzwingen~\cite{ssldragon}. An dieser Stelle greift der
clientseitige Sicherheitsmechanismus: Moderne OpenSSL-Bibliotheken blockieren
standardmäßig die Kommunikation über Protokolle unterhalb von TLS~1.2~\cite{openssl}.
Der Handshake bricht unmittelbar ab. Das Terminal des Clients dokumentiert diesen fatalen
Protokollkonflikt wie in Anhang~\ref{app:sz1-log} aufgezeigt.

In der parallelen Wireshark-Aufzeichnung zeigt sich dieser Abbruch durch ein rot
markiertes Paket mit der Kennzeichnung
\texttt{Alert (Level: Fatal, Description: Protocol Version (70))}.
Der Fehlercode~70 (\textit{protocol\_version}) belegt, dass der Client die erzwungene
Abwärtskompatibilität des Servers zum Schutz des Anwenders blockiert
hat~\cite{alertcode70}.

% ── Szenario 2 ────────────────────────────────────────────────────────────────
\subsection{Szenario~2: Erfolgreicher Downgrade (Schwächung der Mechanismen)}

Das zweite Experiment simuliert den Fall, in dem die clientseitigen Schutzmechanismen
bewusst deaktiviert werden -- beispielsweise, um die Kommunikation mit ungenügend
gehärteten Altsystemen (Legacy-Systemen) im Firmennetzwerk zu
erzwingen~\cite{sentinelone}.

\subsection*{Verbindungsaufbau unter reduzierten Sicherheitsbarrieren}

Nun werden der Server und der Client durch das explizite Herabsetzen des
Sicherheitslevels auf die Stufe~Null (\texttt{@SECLEVEL=0}) jeweils gezwungen, ihre
Schutzschilde zu deaktivieren~\cite{openssl}~(Anhang~\ref{app:sz2-befehle}).

\subsection*{Analyse der Handshake-Schritte}

Durch die Modifikation signalisiert der Client bereits im \textit{Client Hello}, dass er
im Notfall auch das unsichere TLS~1.0 akzeptiert~\cite{ssldragon}. Der Server antwortet
im \textit{Server Hello} mit der Bestätigung der veralteten Version. Da der Client das
nicht mehr blockiert, wechselt die Verbindung in den Zustand \texttt{CONNECTED}.
Anhang~\ref{app:sz2-log} zeigt das Protokoll-Log des Terminals des Clients.

\subsection*{Bewertung der Schwachstelle}

Mit abgeschwächten Sicherheitsregeln ist die Verbindung zwar aufgebaut (\texttt{CONNECTED}), aber nicht mehr
zeitgemäß geschützt. TLS~1.0 bietet keine ausreichende Vertraulichkeit und Integrität,
so dass ein Man-in-the-Middle-Angriff möglich wird~\cite{sentinelone,mdndocs}.

% ── Vergleichstabelle ─────────────────────────────────────────────────────────
\subsection{Analytische Gegenüberstellung der Simulationsergebnisse}

Zur übersichtlichen Strukturierung der Ergebnisse sind die Auswirkungen der beiden Konfigurationen in Tabelle~\ref{tab:vergleich} im Anhang (Abschnitt~\ref{app:vergleichstabelle}) zusammenfassend gegenübergestellt.
